\chapter{Literature Review}\label{chap:background}


\section{Emotion Recognition in Education}


\subsection{Facial Emotion Recognition}
Facial emotion recognition (FER) in education involves the automatic detection of students’ affective and cognitive-emotional states from facial images or video recorded during learning activities. These states include engagement, confusion, boredom, and attention. Recent systematic reviews indicate that FER has become a key artificial intelligence method for emotion assessment in both face-to-face and virtual classrooms, particularly in online and STEM learning environments \cite{anwar2023emotions,vistorte2024integrating}.

Between 2015 and 2017, facial emotion recognition systems in e-learning were mainly based on handcrafted feature descriptors and traditional classifiers. Early studies applied supervised learning methods such as SVM, kNN, random forests, and RBF neural networks to classify basic emotions and moods, including happiness, sadness, and anger, and to relate them to students’ learning progress in e-learning platforms \cite{ayvaz2017facial,ashwin2015elearning,alalwani2016mood}. These systems generally relied on generic facial expression datasets and focused on basic emotion labels rather than fine-grained academic emotions. Therefore, facial emotion was mainly treated as a supportive signal for personalization and for tracking students’ general mood. In parallel, early CNN-based facial expression recognition methods were introduced to address limited training data and sample-order effects, providing the initial deep learning foundation later adopted in educational contexts \cite{lopes2017facial}.

From 2019 to 2020, facial emotion recognition in education shifted toward deeper architectures and more educationally relevant targets. Classroom-based studies employed CNNs and hybrid CNN/Bi-LSTM models to identify students’ affective states, including positive and negative classroom climate, engagement, boredom, and inattentiveness, directly from lecture videos \cite{ramakrishnan2019toward,ashwin2020automatic}. These approaches typically combined CNN-based spatial feature extraction with recurrent layers to capture temporal dynamics. Although these methods still relied largely on general FER datasets and basic emotion categories, they marked a clear transition from general facial expression analysis toward educationally relevant constructs and classroom observation \cite{anwar2023emotions}.

Between 2021 and 2022, deep learning for facial emotion recognition became more closely integrated with distance education environments. CNN-based FER was established as the dominant approach in computer vision, with end-to-end models trained on benchmark datasets such as FER2013, CK+, and AffectNet \cite{liu2023application}. At the same time, research on children’s emotion recognition using multiple deep learning models and explainable AI emphasized the importance of interpretable attention maps and age-appropriate modelling in educational and related contexts \cite{rathod2022kids}. Domain-specific FER for learning settings was also explored by Akputu et al. \cite{akputu2022recognizing}, who combined handcrafted descriptors with multiple kernel learning and evaluated their approach on both general datasets and a simulated learning-emotion dataset designed for educational scenarios. Overall, these studies reflected a clearer shift toward learning-specific affect recognition while also stressing the importance of interpretability and contextual relevance.

From 2023 onward, facial emotion recognition has been used more widely in online learning and real-time engagement detection. In online lectures, Siswantoro et al. \cite{siswantoro2024facial} used a CNN-based model to classify students as engaged or disengaged in real time, helping instructors monitor attention and participation continuously. Zhao et al. \cite{zhao2025context} further developed this area by introducing a context-aware academic emotion dataset and benchmark designed specifically for learning-related emotions, showing that models trained on such data perform better than those trained only on general FER datasets. In addition, Salloum et al. \cite{salloum2025emotion} proposed an improved CNN model based on the FER2013 dataset to classify seven basic emotions and discussed its use in AI systems that can adjust teaching methods according to students’ detected emotions. Although these studies show the value of FER as a non-intrusive and real-time tool, challenges still remain, including lighting conditions, occlusion, head pose, cultural differences, privacy issues, and the continued use of general FER datasets in many systems \cite{anwar2023emotions,vistorte2024integrating}. These developments have encouraged a shift from facial analysis alone toward multimodal emotion recognition in educational settings.

\section{Adaptive Learning Systems}

Adaptive learning in education involves the use of computational systems that personalize instructional content, feedback, and pacing according to individual learner characteristics, such as prior knowledge, engagement, affective states, and learning preferences. This paradigm emerged in the early 2000s in response to the challenge of learner heterogeneity, including differences in learning pace, background knowledge, motivation, and support needs that traditional instruction could not effectively address \cite{Strielkowski2024AI‐driven,Gligorea2023Adaptive}. Its central goal is to enhance educational outcomes by delivering a personalized learning experience through real-time adaptation mechanisms.


In the foundational period, from the late 1990s to the early 2010s, adaptive learning systems were mainly built on rule-based engines that encoded expert knowledge into explicit decision rules. These systems relied on structured learner models, which represented factors such as prior achievement and preferred learning style, to guide adaptation. For instance, many early intelligent tutoring systems selected instructional actions based on students’ responses to quizzes and exercises \cite{barbosa2023adaptive}. Fuzzy logic was also introduced to address uncertainty in estimating learner states, allowing more flexible classification into categories such as “beginner” or “advanced” rather than using strict binary labels \cite{Strielkowski2024AI‐driven}. At this stage, common input data included quiz scores, clickstream data, and self-reported preferences, while adaptation mainly focused on adjusting content difficulty and sequencing according to these cognitive indicators.

From the mid-2010s to 2023, adaptive learning experienced a major transformation through the integration of artificial intelligence techniques. Systems increasingly used machine learning methods, such as clustering to group similar learners and collaborative filtering to recommend appropriate learning resources, enabling more dynamic personalization \cite{Essa2023Personalized}. Neural networks were also applied to automatically identify learning styles from behavioral attributes \cite{Li2025Intelligent}. At the same time, learning analytics became central to adaptive learning, as platforms collected extensive interaction data, including time on task, assessment results over time, and engagement indicators, to continuously update learner models. This development allowed adaptive systems to modify content sequencing and difficulty in real time according to students’ ongoing performance. Commercial platforms such as Knewton and Smart Sparrow exemplified this shift by using large scale educational data to provide personalized resource recommendations \cite{Jing2023Research}. Furthermore, integration with Learning Management Systems (LMS) supported scalable implementation across a wide range of educational contexts.


From around 2022 to 2025, adaptive learning systems increasingly moved beyond purely cognitive adaptation by incorporating affective computing and multimodal learner-state analysis. In this stage, systems combined traditional behavioral data, such as quiz performance, time-on-task, and interaction logs, with affect-related inputs including facial expressions from webcams, sentiment from chat or text interactions, and, in some cases, physiological signals such as heart rate or skin conductance \cite{Akavova2023Adaptive,Rane2023Education,Gkintoni2025Challenging}. These data streams were processed through multimodal fusion modules to detect states such as frustration, disengagement, or low attention in real time, allowing intelligent tutoring systems to adapt feedback, pacing, or content difficulty more responsively. This marked an important transition from systems that adapted only to knowledge and performance toward systems that also responded to learners’ emotional conditions.

From 2024 onward, adaptive learning has been further transformed by generative AI, large language models, agent-based planning, and neuro-symbolic methods. LLM-based tutors can now support multi-turn personalized dialogue, while agent-based systems dynamically plan instructional strategies such as scaffolding, prompting, or Socratic questioning according to the learner’s progress \cite{Kabudi2021AI-enabled,Khosravi2020Development,Katsaris2021Adaptive}. Some recent systems also integrate conversational history, performance data, and physiological or affective signals to generate more personalized support, while neuro-symbolic architectures combine neural models with symbolic reasoning or knowledge graphs to improve explainability in learner modelling and instructional decisions \cite{Khosravi2020Development,Katsaris2021Adaptive}.
\section{challenges in emotion recognation for adaptive learning }
Despite recent advances in emotion-adaptive educational technologies, current systems remain limited in their ability to integrate affective signals into real-time instructional decision-making. In most educational applications, facial expression recognition is the dominant approach, sometimes combined with physiological or behavioral cues, to monitor learners’ states and support broad adjustments in feedback or instructional interaction \cite{Hua2025Multimodal}. By comparison, text-based and discourse-derived emotion detection has been explored mainly in wider natural language processing and human-computer interaction contexts for analytic and monitoring purposes, with relatively limited use in real-time pedagogical adaptation within learning environments \cite{Acheampong2020Text‐based,Kusal2023A}. Furthermore, multimodal approaches that jointly combine facial and text-based affective indicators within a principled and fine-grained adaptive framework—one that can dynamically regulate pacing, difficulty, feedback, and hints in authentic learning settings—remain largely underexplored. Consequently, there is still a strong need for rigorously evaluated emotion-aware adaptive learning systems that better exploit the complementary value of face and text modalities in real time to deliver pedagogically meaningful support \cite{vistorte2024integrating}.